\section{Lexical Multi-Dimensional Analysis Overview}

Lexical multidimensional analysis (LMDA), an extension of the multidimensional (MD) analysis framework developed by Biber in the 1980s (``multi-feature multidimensional analysis'') to study register variation. Through the identification of (lexical) dimensions or sets of correlated lexical features, LMDA enables the analysis of lexical patterning from a multidimensional perspective. These lexical dimensions represent a variety of latent, macro-level discursive constructs. \parencite[p.1]{berbersardinhaLexicalMultiDimensionalAnalysis2025}

A typical LMDA comprises the procedures described in Figure \ref{fig:lexical_md_analysis_procedures}.

\import{diagrams/}{md_analysis_diagram_lexical_procedures}

\section{Lexical Variable Selection via Keywords}

Figure \ref{fig:lexical_variables_selection_by_keywords} describes the selection of lexical variables.

\import{diagrams/}{lexical_variables_selection_by_keywords}

This process is implemented with two programmes in the LMDA pipeline:

\begin{itemize}
    \item \textit{keylemmas.py}
    \item \textit{select\_kws\_stratified.py}
\end{itemize}

\subsection{\textit{keylemmas.py}}

\textit{keylemmas.py} computes key lemmas from the tagged corpus.

The programme compares each subcorpus against all other subcorpora combined and produces a ranked key-lemma table for each subcorpus.

The output is used by \textit{select\_kws\_stratified.py} to choose a balanced list of positive keywords across subcorpus strata.

Click \href{https://github.com/laelgelc/cl_st1_anna/blob/main/cl_st1_ph2_anna/cl_st1_ph2_anna_pipeline.md#3-extract-key-lemmas-by-group}{\textit{Extract key lemmas by group}} for an example.

\subsection{\textit{select\_kws\_stratified.py}}

\textit{select\_kws\_stratified.py} selects a balanced set of positive keywords from subcorpus-specific key-lemma tables.

It reads the key-lemma tables produced by \textit{keylemmas.py}, extracts lemmas classified as positive keywords (`POSKW'), applies filtering rules, enforces an equal maximum quota per subcorpus, and writes:

\begin{enumerate}
    \item one selected keyword list per subcorpus
    \item one consolidated deduplicated keyword list for downstream LMDA processing
\end{enumerate}

Click \href{https://github.com/laelgelc/cl_st1_anna/blob/main/cl_st1_ph2_anna/cl_st1_ph2_anna_pipeline.md#4-select-a-stratified-keyword-set}{\textit{Select a stratified keyword set}} for an example.
